\documentclass[]{article}
\usepackage[a4paper, total={15cm,23cm}]{geometry}
\usepackage{fancyhdr}
\usepackage{graphicx}
\usepackage{amsmath}
\usepackage{amssymb}
\usepackage{xcolor}
\usepackage{tikz}
\usepackage{verbatim}
\usepackage{tcolorbox}
\usepackage{textcomp}
\usepackage{xcomment}
\usepackage{xstring}
\usepackage{array}
%opening
\title{Motion Diagram Jogging}
\author{Benjamin Bauml}
\date{Spring 2024}
\pagestyle{fancy}
\rhead{PH 221}
\chead{Spring 2024}
\lhead{Week 1}

% Version 2024-02-21
% Changes
% 2024-02-21 Added xstring package to enable smooth implementation of new \ModePage command.
% For Assignment, leave Purpose as 1. For Worksheet, set to 2. For Student Solution, set to 3. For Teacher Solution, set to 4.
\newcommand{\Purpose}{4}

\newcommand{\Exclusion}{0}
\newcommand{\PageTurn}{0}
\newcommand{\GrayProb}{0}
\newcommand{\Tipsy}{0}

% Assignment
\if\Purpose1
\renewcommand{\Exclusion}{1}
\fi
% Worksheet
\if\Purpose2
\renewcommand{\Exclusion}{1}
\renewcommand{\PageTurn}{1}
\fi
% Student Solution
\if\Purpose3
\renewcommand{\PageTurn}{1}
\renewcommand{\GrayProb}{1}
\fi
% Teaching Copy
\if\Purpose4
\renewcommand{\PageTurn}{1}
\renewcommand{\GrayProb}{1}
\renewcommand{\Tipsy}{1}
\fi

\if\Exclusion1
\xcomment{Title,Problem,ProblemSub,PassFig}
\fi

\def \NewQ {0}
\def \PForce {0}
\newcommand{\MaybePage}[1]{
	\def \PForce {#1}
	\if\PForce1
		\newpage
	\else
		\if\NewQ0
		\gdef \NewQ {\PageTurn}
		\else
		\newpage
		\fi
	\fi
}

\newcommand{\ModePage}[1]{
	\IfSubStr{#1}{\Purpose}{\newpage}{}
}

\newenvironment{Problem}[2][0]{%The first argument is optional, and if it is set to 1, the \newpage will be forced.
\MaybePage{#1}
\section*{#2}
\if\GrayProb1
\begin{tcolorbox}[colback=lightgray,colframe=lightgray,sharp corners,boxsep=1pt,left=0pt,right=0pt,top=0pt,bottom=0pt,after skip=2pt]
\else
\begin{tcolorbox}[colback=white,colframe=white,sharp corners,boxsep=1pt,left=0pt,right=0pt,top=0pt,bottom=0pt,after skip=2pt]
\fi
}{
\end{tcolorbox}\noindent
}

\newenvironment{ProblemSub}[1][0]{%The argument is optional, and if a string of numbers is entered into it, it will force a \newpage in any \Purpose that shows up in the string. For example, "13" would lead to the newpage being forced in modes 1 and 3.
\ModePage{#1}
\if\GrayProb1
\begin{tcolorbox}[colback=lightgray,colframe=lightgray,sharp corners,boxsep=1pt,left=0pt,right=0pt,top=0pt,bottom=0pt,after skip=2pt]
\else
\begin{tcolorbox}[colback=white,colframe=white,sharp corners,boxsep=1pt,left=0pt,right=0pt,top=0pt,bottom=0pt,after skip=2pt]
\fi
}{
\end{tcolorbox}\noindent
}

\newenvironment{PassFig}{\begin{figure}[h]}{\end{figure}}

\newcommand{\TeachingTips}[1]{
\if\Tipsy1
\begin{tcolorbox}[colback=lightgray,colframe=black]
#1
\end{tcolorbox}
\fi
}

\newenvironment{Title}{\maketitle}{}

\begin{document}
\begin{Title}
\begin{center}
	This material is borrowed/adapted from PH 201 Tutorial 1 for Fall 2020.
\end{center}
\end{Title}

\begin{Problem}{Activity}
	Mary runs east at 6.0 km/h for 0.50 hr. She then turns around and runs west at 4.0 km/h for 1.0 h. Let east be the positive direction.
\end{Problem}
\begin{ProblemSub}
	a) Draw a motion diagram for Mary. Assume turning around takes negligible time.
\end{ProblemSub}
\begin{figure}[h]
	\centering
	\begin{tikzpicture}[scale=1.5]
		\node[anchor=south] at (0,0) {\phantom{,}0\phantom{,}};
		\filldraw[black] (0,0) circle (2pt);
		\node[anchor=south] at (1,0) {1,6};
		\filldraw[black] (1,0) circle (2pt);
		\node[anchor=south] at (2,0) {\phantom{,}2\phantom{,}};
		\filldraw[black] (2,0) circle (2pt);
		\node[anchor=south] at (3,0) {\phantom{,}3\phantom{,}};
		\filldraw[black] (3,0) circle (2pt);
		\node[anchor=south] at (2.33,0) {\phantom{,}4\phantom{,}};
		\filldraw[black] (2.33,0) circle (2pt);
		\node[anchor=south] at (1.66,0) {\phantom{,}5\phantom{,}};
		\filldraw[black] (1.66,0) circle (2pt);
		%\node[anchor=south] at (1,0) {6};
		%\filldraw[black] (1,0) circle (2pt);
		\node[anchor=south] at (0.33,0) {\phantom{,}7\phantom{,}};
		\filldraw[black] (0.33,0) circle (2pt);
		\node[anchor=south] at (-0.33,0) {\phantom{,}8\phantom{,}};
		\filldraw[black] (-0.33,0) circle (2pt);
		\node[anchor=south] at (-1,0) {\phantom{,}9\phantom{,}};
		\filldraw[black] (-1,0) circle (2pt);
	\end{tikzpicture}
\end{figure}

Here, we have chosen $ \Delta t = 10 $ min. Other choices are also acceptable.
\begin{ProblemSub}
	b) How far does Mary run?
\end{ProblemSub}
To find how far she runs, we want the total distance she ran. Distance doesn't include directional information, so we just find how far she ran east and how far she ran west (two positive numbers) and add them together:
\[
\begin{split}
	d_{E} & = (6.0 \text{ km/h})(0.50\text{ h}) = 3.0\text{ km} \\
	d_{W} & = (4.0 \text{ km/h})(1.0\text{ h}) = 4.0\text{ km} \\
	d & = d_{E} + d_{W} = 3.0\text{ km} + 4.0\text{ km} = 7.0\text{ km.}
\end{split}
\]
\begin{ProblemSub}
	c) What is Mary’s displacement?
\end{ProblemSub}
Displacement is a vector, so we need to add the displacement of her eastward run (a positive number) to the displacement of her westward run (a negative number) to get the total displacement:
\[
\begin{split}
	\vec{d}_{E} & = (6.0 \text{ km/h})(0.50\text{ h}) = 3.0\text{ km} \\
	\vec{d}_{W} & = (-4.0 \text{ km/h})(1.0\text{ h}) = -4.0\text{ km} \\
	\vec{d} & = \vec{d}_{E} + \vec{d}_{W} = 3.0\text{ km} + (-4.0\text{ km}) = -1.0\text{ km}, \text{ i.e. } 1.0\text{ km, west.}
\end{split}
\]
\end{document}